Se o mecanismo dominante fosse \textbf{troca térmica com o ambiente}, a escala principal seria de \textbf{área superficial}, não de volume, pois o calor é transferido pela superfície.

Com fator linear $L_{\rm G}/L_{\rm L} = 12$:
\[
A \sim L^2 \implies \frac{A_{\rm G}}{A_{\rm L}} = 12^2 = 144, \qquad V \sim L^3 \implies \frac{V_{\rm G}}{V_{\rm L}} = 12^3 = 1728.
\]

\textbf{Comparação das conclusões:}
\begin{itemize}
\item Hipótese original (metabolismo $\propto$ massa/volume): demanda $\times 1728$.
\item Hipótese térmica (troca de calor $\propto$ área superficial): demanda $\times 144$.
\item Razão entre as duas estimativas: $1728/144 = 12$.
\end{itemize}

Portanto, se a demanda energética de Gulliver fosse ditada pela perda de calor para o ambiente, os liliputianos teriam concluído que ele precisa de apenas $\mathbf{144}$ vezes mais alimento, não 1728 vezes. A conclusão é dramaticamente diferente, ilustrando como a hipótese de escala adotada (volume versus área) altera o resultado.

